\exercisesection{Algebras of Continuous Functions on a Compact Space}

All Banach spaces and all Banach algebras below are complex.

\begin{enumerate}
\def\labelenumi{\arabic{enumi})}
\tightlist
\item
  Let X be a topological space, let \(\mathcal{C}_{b}(\mathrm{X})\) be the stellar algebra of bounded continuous complex functions on X, and let \(\widehat{\mathrm{X}} = \mathrm{X}(\mathcal{C}_{b}(\mathrm{X}))\).
\end{enumerate}

\begin{enumerate}
\def\labelenumi{\alph{enumi})}
\tightlist
\item
  Let A be a unital Banach subalgebra of \(\mathcal{C}_{b}(\mathrm{X})\) and let \(R_{A}\) be the equivalence relation on X such that x is equivalent to \(x'\) if \(f(x) = f(x')\) for every function \(f \in A\). The \(R_{A}\)-saturation S of a compact subset K of X is closed. (Let \(x \in \overline{S}\); for all \(f \in A\) and every \(\varepsilon > 0\), let \(V_{f,\varepsilon}\) be the set of \(x' \in X\) such that \(|f(x) - f(x')| \leqslant \varepsilon\); if \(f_{1}, \ldots, f_{n} \in A\) and \(\varepsilon_{1}, \ldots, \varepsilon_{n} > 0\), one has
\end{enumerate}

\[
\mathrm{V} _ {f _ {1}, \varepsilon_ {1}} \cap \dots \cap \mathrm{V} _ {f _ {n}, \varepsilon_ {n}} \cap \mathrm{K} \neq \varnothing ,
\]

hence

\[
\bigcap_{\substack{f\in \mathrm{A}\\ \varepsilon >0}}\mathrm{V}_{f,\varepsilon}\cap \mathrm{K}\neq \varnothing
\]

whence \(x \in S\).)

\begin{enumerate}
\def\labelenumi{\alph{enumi})}
\setcounter{enumi}{1}
\item
  Assume X is normal. Let R be a closed equivalence relation on X and let \(A = A_{R}\) be the subalgebra of \(\mathcal{C}_{b}(X)\) formed by the functions constant on the classes modulo R. Show that \(R = R_{A}\) (Use TG, IX, p.~105, exerc. 19, a).
\item
  Assume X is compact. \(R \mapsto A_{R}\) is a bijection from the set of separated equivalence relations on X onto the set of stellar subalgebras of \(\mathcal{C}(X)\) (Let A be a stellar subalgebra of \(\mathcal{C}(X)\). Show that \(R_{A}\) is separated, and apply the Weierstrass–Stone theorem to the algebra of continuous functions on \(X/R_{A}\) obtained from A by passing to the quotient.)
\end{enumerate}

\begin{enumerate}
\def\labelenumi{\arabic{enumi})}
\setcounter{enumi}{1}
\tightlist
\item
  Let X be a compact topological space, R a separated equivalence relation on X, and B a closed subalgebra of \(\mathcal{C}(\mathrm{X})\) containing \(\mathrm{A}_{\mathrm{R}}\) (notation of Exercise 1). If \(f \in \mathcal{C}(\mathrm{X})\) coincides on each R-class \(c\) with an element \(g_{c}\) of B, then \(f \in \mathrm{B}\). (Let \(\varepsilon > 0\). There exists a saturated open neighborhood \(\mathrm{V}_{c}\) of \(c\) such that \(|f - g_{c}| \leqslant \varepsilon\) on \(\mathrm{V}_{c}\). Let \(\mathrm{V}_{c_{1}}, \ldots, \mathrm{V}_{c_{n}}\) be a cover of X. Let \((u_{1}, \ldots, u_{n})\) be a partition of unity subordinate to \((\mathrm{V}_{c_{1}}, \ldots, \mathrm{V}_{c_{n}})\) and consisting of functions of \(\mathrm{A}_{\mathrm{R}}\). Then
\end{enumerate}

\[
g = u _ {1} g _ {c _ {1}} + \dots + u _ {n} g _ {c _ {n}} \in \mathrm{B},
\]

and \(|f - g| \leqslant \varepsilon\).)

\begin{enumerate}
\def\labelenumi{\arabic{enumi})}
\setcounter{enumi}{2}
\tightlist
\item
  Let X be a compact topological space and, for every \(x \in X\), let \(A_{x}\) be a commutative unital Banach algebra, with \(e_{x}\) denoting the identity element. We assume \(\|e_{x}\| = 1\). Let B be the product Banach algebra of the \(A_{x}\) for \(x \in X\). Let C be a Banach subalgebra of B having the following properties:
\end{enumerate}

\begin{enumerate}
\def\labelenumi{(\roman{enumi})}
\item
  \(e = (e_x)_{x \in \mathbf{X}} \in \mathbf{C};\)
\item
  If \(a = (a_x) \in \mathrm{C}\) and \(f \in \mathcal{C}(\mathrm{X})\) , the function \(fa: x \mapsto f(x)a_x\) is an element of \(\mathrm{C}\) ;
\item
  For every \(a = (a_x) \in \mathrm{C}\), the function \(x \mapsto \|a_x\|\) is upper semicontinuous.
\end{enumerate}

Let \(\widehat{\mathbf{X}}\) be the disjoint union of the \(\mathsf{X}(\mathsf{A}_x)\). For every \(x_0 \in \mathbf{X}\) and \(\chi \in \mathsf{X}(\mathsf{A}_{x_0})\), let \(\varphi(\chi)\) be the character of C defined by \(\varphi(\chi)((a_x)_{x \in \mathbf{X}}) = \chi(a_{x_0})\). Prove that \(\varphi\) defines a bijection from \(\widehat{\mathbf{X}}\) onto \(\mathsf{X}(\mathbf{C})\). (Let \(\xi \in \mathsf{X}(\mathbf{C})\). Since \(\mathcal{C}(\mathbf{X})\) can be identified with a closed subalgebra of C, \(\xi\) defines a character of \(\mathcal{C}(\mathbf{X})\), hence there exists \(x_0 \in \mathbf{X}\). Let \(a = (a_x) \in \mathbf{C}\) and \(b = (b_x) \in \mathbf{C}\), with \(a_{x_0} = b_{x_0}\); let \(\varepsilon > 0\); we have \(\|a_x - b_x\| < \varepsilon\) on a neighborhood U of \(x_0\); let \(f \in \mathcal{C}(\mathbf{X})\) be such that \(0 \leqslant f \leqslant 1\), equal to 1 at \(x_0\) and to 0 outside U; we have \(\|fa - fb\| \leqslant \varepsilon\), whence \(|\xi(fa - fb)| \leqslant \varepsilon\), whence \(|\xi(a) - \xi(b)| \leqslant \varepsilon\), whence \(\xi(a) = \xi(b)\). Deduce that there exists a \(\chi \in \mathsf{X}(\mathsf{A}_{x_0})\) such that \(\xi = \varphi(\chi)\).

\begin{enumerate}
\def\labelenumi{\arabic{enumi})}
\setcounter{enumi}{3}
\tightlist
\item
  Let X be a compact topological space such that the Banach algebra \(\mathcal{C}(\mathrm{X})\) is generated by a sequence \((j_n)\) of idempotents. The function
\end{enumerate}

\[
\omega \mapsto \sum_ {n = 1} ^ {\infty} \frac {1}{3 ^ {n}} (2 j _ {n} (\omega) - 1)
\]

separates the points of X. It generates the Banach algebra \(\mathcal{C}(\mathrm{X})\).

\begin{enumerate}
\def\labelenumi{\arabic{enumi})}
\setcounter{enumi}{4}
\tightlist
\item
  Let A be the Banach algebra of functions from {[}0,1{]} to C admitting continuous derivatives on {[}0,1{]} up to order n (example 4 of I, p.~18). Let \(\omega \in [0,1] = \mathsf{X}(\mathsf{A})\).
\end{enumerate}

\begin{enumerate}
\def\labelenumi{\alph{enumi})}
\item
  The smallest closed ideal I of A such that \(V(I) = \{\omega\}\) is the set of \(f \in A\) such that \(f, f', \ldots, f^{(n)}\) vanish at \(\omega\). (Use prop. 5 of I, p.~94.)
\item
  If \(g \in A\), the norm of its canonical image in A/I is equal to
\end{enumerate}

\[
\sum_ {k = 0} ^ {n} \frac {| g ^ {(k)} (\omega) |}{k !}.
\]

\begin{enumerate}
\def\labelenumi{\alph{enumi})}
\setcounter{enumi}{2}
\tightlist
\item
  The algebra A/I is isomorphic to \(\mathbf{C}[\mathrm{X}]/(\mathrm{X}^{n+1})\); the only maximal ideal of A containing I is the set of \(f \in A\) such that \(f(\omega) = 0\).
\end{enumerate}

\begin{enumerate}
\def\labelenumi{\arabic{enumi})}
\setcounter{enumi}{5}
\tightlist
\item
  Let \(\Delta\) be the disk \(|\zeta| \leqslant 1\) in \(\mathbf{C}\) and A the Banach algebra of continuous complex-valued functions on \(\Delta\) that are analytic in the interior of \(\Delta\) (example 9 of I, p.~20).
\end{enumerate}

\begin{enumerate}
\def\labelenumi{\alph{enumi})}
\tightlist
\item
  If \(x \in A\), then \(x\) is the limit in A of the functions
\end{enumerate}

\[
\zeta \mapsto x _ {n} (\zeta) = x \left(\frac {n \zeta}{n + 1}\right).
\]

The identity mapping of \(\Delta\) is a generator of A.

\begin{enumerate}
\def\labelenumi{\alph{enumi})}
\setcounter{enumi}{1}
\item
  The space X(A) is canonically identified with \(\Delta\) (Use assertion g) of prop. 1 of I, p.~142.)
\item
  Let S be a closed subset of \(\Delta\), identified with a closed subset of \(\mathsf{X}(\mathsf{A})\). If S has a non-isolated point \(\zeta\) such that \(|\zeta| < 1\), the closure of S for the Jacobson topology is \(\Delta\).
\end{enumerate}

\(d)\) For \(x \in \mathrm{A}\), set \(x^{*}(\zeta) = x(\overline{\zeta})\). Then \(x \mapsto x^{*}\) is an isometric involution of A. The Hermitian characters of A are given by the real elements of \(\Delta\).

\begin{enumerate}
\def\labelenumi{\alph{enumi})}
\setcounter{enumi}{4}
\tightlist
\item
  The map \(x \mapsto x|U\) is an isomorphism from A onto \(\mathrm{P}(\mathbf{U})\). Every function in \(\mathcal{C}_{\mathbf{R}}(\mathbf{U})\) is the uniform limit of real parts of functions in \(\mathrm{P}(\mathbf{U})\).
\end{enumerate}

\begin{enumerate}
\def\labelenumi{\arabic{enumi})}
\setcounter{enumi}{6}
\tightlist
\item
  Let \(A_{1}\) (resp. \(A_{2}\)) be the unital Banach algebra considered in Exercise 5 (resp. 6). Let \(A = A_{1} \times A_{2}\). Show that A is semisimple, but that the image of the Gelfand transform of A is neither closed nor dense in \(\mathcal{C}(\mathsf{X}(\mathsf{A}))\).
\end{enumerate}

¶ 8) Let X be a compact topological space and B a unital Banach subalgebra of \(\mathcal{C}(\mathrm{X})\) (with the induced norm), separating the points of X. One identifies X with a closed subset of \(\mathrm{X(B)} = \mathrm{X}'\). Let \(\mathrm{B}^*\) be the set of invertible elements of B. One says that B is logmodular if the set of functions \(\log(|f|)\), for \(f \in \mathrm{B}^*\), is dense in \(\mathcal{C}_{\mathbf{R}}(\mathrm{X})\). (This is the case, in particular, if the set of functions \(\mathcal{R}(f)\), for \(f \in \mathrm{B}\), is dense in \(\mathcal{C}_{\mathbf{R}}(\mathrm{X})\), because \(\log(|e^f|) = \mathcal{R}(f)\).) In what follows, one assumes that B is logmodular.

\begin{enumerate}
\def\labelenumi{\alph{enumi})}
\tightlist
\item
  For every \(\chi \in \mathrm{X}'\) , there exists a unique measure \(\mu_{\chi} \geqslant 0\) on X such that
\end{enumerate}

\[
\log (| \chi (f) |) = \int_ {\mathrm{X}} \log (| f (\omega) |) d \mu_ {\chi} (\omega)
\]

for all \(f \in \mathrm{B}^*\), hence such that \(\chi(f) = \int_{\mathrm{X}} f(\omega) d\mu_{\chi}(\omega)\) for all \(f \in \mathrm{B}\). \(b)\) One has

\[
\log | \chi (f) | \leqslant \int_ {\mathrm{X}} \log | f (\omega) | d \mu_ {\chi} (\omega)
\]

for all \(f \in \mathrm{B}\) and the condition \(\chi(f) = \int_{\mathrm{X}} f(\omega) d\mu_{\chi}(\omega)\) for all \(f \in \mathrm{B}\) defines \(\mu_{\chi}\) uniquely.

\begin{enumerate}
\def\labelenumi{\alph{enumi})}
\setcounter{enumi}{2}
\tightlist
\item
  Let \(\chi \in X'\) and \(\mu\) be a positive measure on \(X\). Write \(\mu = \mu_1 + \mu_2\), where \(\mu_1\) is absolutely continuous with respect to \(\mu_{\chi}\), and where \(\mu_2\) and \(\mu_{\chi}\) are mutually singular. Let I be the kernel of \(\chi\). Then
\end{enumerate}

\[
\inf _ {f \in \mathrm{I}} \int_ {\mathrm{X}} | 1 - f | ^ {2} d \mu = \inf _ {f \in \mathrm{I}} \int_ {\mathrm{X}} | 1 - f | ^ {2} d \mu_ {1}.
\]

\(d)\) Let \(\chi\in\mathrm{X}^{\prime}\) and let \(h\) be a nonnegative function in \(\mathcal{L}^{1}(\mu_{\chi})\). Let I be the kernel of \(\chi\). Then

\[
\inf _ {f \in \mathrm{I}} \int_ {\mathrm{X}} | 1 - f | ^ {2} h d \mu_ {\chi} = \exp \Bigl (\int_ {\mathrm{X}} \log (h) d \mu_ {\chi} \Bigr).
\]

\begin{enumerate}
\def\labelenumi{\alph{enumi})}
\setcounter{enumi}{4}
\tightlist
\item
  For every positive measure \(\mu\) on X, and every \(p \in [1, +\infty[\), let \(\mathrm{H}^p(\mu)\) be the closure of the canonical image of B in \(\mathrm{L}^p(\mu)\). Let \(\chi \in \mathrm{X}'\) and let I be the kernel of \(\chi\). Then \(\mathrm{L}^2(\mu_\chi)\) is the Hilbert sum of \(\mathrm{H}^2(\mu_\chi)\) and of the closure in \(\mathrm{L}^2(\mu_\chi)\) of the canonical image of \(\bar{\mathrm{I}}\) (the set of conjugates of functions belonging to I).
\end{enumerate}

\(f)\) The space \(\mathrm{H}^1 (\mu_\chi)\) is the set of classes \(\widetilde{h}\) of \(h\in \mathcal{L}^1 (\mu_\chi)\) such that \(\int_{\mathrm{X}}fh  d\mu_{\chi} = 0\) for all \(f\in \mathrm{I}\).

\(g)\) Let \(h\) be a nonnegative function in \(\mathcal{L}^1 (\mu_\chi)\). Then we have \(\widetilde{h} = |\widetilde{f}|\) with \(\widetilde{f}\in \mathrm{H}^1 (\mu_\chi)\) and \(\int_{\mathrm{X}}f  d\mu_{\chi}\neq 0\) if and only if \(\log (h)\in \mathcal{L}^1 (\mu_\chi)\).

¶ 9) Let X be a compact topological space and B a unital Banach subalgebra of \(\mathcal{C}(\mathrm{X})\) (with the induced norm). Every element \(x\) of X defines a character \(\chi_x\) of B. If \(x\) and \(x'\) are in X, we write \(x \sim x'\) when \(\| \chi_x - \chi_{x'} \| \neq 2\). a) Let \(x\) and \(x'\) be in X. If there exists a sequence \((f_n)\) of elements of B such that \(\| f_n \| \leqslant 1\), \(|f_n(x)|\) tends to 1, and \(\liminf_{n \to \infty} |f_n(x) - f_n(x')| > 0\), then \(x\not\sim x'\). (One may assume that \(f_n(x)\) tends to 1. Consider \(g_n = (f_n - \lambda_n)(1 - \lambda_n f_n)^{-1}\), where \((\lambda_n)\) is a sequence of numbers in \([0,1[\) tending toward 1. We have \(\| g_n \| \leqslant 1\). If the sequence \((\lambda_n)\) is well chosen, \(g_n(x)\) tends to 1 and \(g_n(x')\) tends to -1.)

\begin{enumerate}
\def\labelenumi{\alph{enumi})}
\setcounter{enumi}{1}
\item
  Let R be the set of real parts of the elements of B. If \(x \sim x'\), there exists \(c \in ]0,1[\) such that \(f(x) \geqslant cf(x')\) and \(f(x') \geqslant cf(x)\) for every \(f \geqslant 0\) in R. (It suffices to obtain the first condition; if it is not true, there exists a sequence \((f_n)\) in R with \(f_n \geqslant 0\), \(f_n(x') = 1\), and \(f_n(x)\) tending to 0; let \(g_n \in B\) with \(\mathcal{R}(g_n) = f_n\). Then \(e^{-g_n} \in \mathrm{B}\), \(\| e^{-g_n} \| \leqslant 1\), \(|e^{-g_n(x')}| = 1 / e\), and \(|e^{-g_n(x)}|\) tends to 1, which contradicts a).)
\item
  There exist positive measures \(\alpha\) and \(\beta\) on X such that \(f(x)-cf(x')=\alpha(f)\) and \(f(x')-cf(x)=\beta(f)\) for \(f\in R\) . Setting \(\mu_{x}=(1-c^{2})^{-1}(c\beta+\alpha)\) , \(\mu_{x'}=(1-c^{2})^{-1}(c\alpha+\beta)\) , one has \(\mu_{x}(f)=f(x)\) and \(\mu_{x'}(f)=f(x')\) for all \(f\in B\) , and \(c\mu_{x}\leqslant\mu_{x'}\) , \(c\mu_{x'}\leqslant\mu_{x}\) .
\item
  We still assume \(x \sim x'\). Show that if \(\mu\) is a positive measure on X such that \(\mu(f) = f(x)\) for all \(f \in B\), there exists a positive measure \(\mu'\) on X such that \(\mu'(f) = f(x')\) for all \(f \in B\) and \(c\mu \leqslant \mu'\) for some \(c \in]0,1[\). (With the notations of c), set \(\mu' = \mu_{x'} - c\mu_x + c\mu\).)
\item
  The relation \(x \sim x'\) is an equivalence relation on X. (Use d).)
\end{enumerate}

¶ 10) Let X be a compact topological space and A a unital Banach subalgebra of \(\mathcal{C}(\mathrm{X})\) for the induced norm.

\begin{enumerate}
\def\labelenumi{\alph{enumi})}
\item
  A subset K of X is said to be antisymmetric (relative to A) if every function \(f \in A\) which is real on K is constant on K. Every antisymmetric subset is contained in a maximal antisymmetric subset. A maximal antisymmetric subset is closed. The set \(\mathfrak{R}\) of maximal antisymmetric subsets is a partition of X.
\item
  Let \(A^{\perp}\) be the orthogonal of A in the Banach space of complex measures on X. Let \(\mu\) be an extreme point of the unit ball of \(A^{\perp}\). Show that the support of \(\mu\) is antisymmetric.
\item
  Let \(f \in \mathcal{C}(\mathrm{X})\). If \(f|K \in A|K\) for all \(K \in \mathfrak{R}\), then \(f \in A\). (Apply b) and the Krein–Milman theorem.) From this, recover the Weierstrass–Stone theorem (TG, X, p.~36, th. 3).
\end{enumerate}

\(d)\) A nonempty subset E of X is called a peak (relative to A) if there exists \(f\in \mathrm{A}\) such that \(\| f\| = 1\) and E is the set of those \(x\in \mathrm{X}\) for which \(f(x) = 1\). One may then assume \(|f(x)| < 1\) for \(x\notin \mathrm{E}\) (replace \(f\) by \(\frac{1}{2} (1 + f)\)). Every nonempty countable intersection of peaks is a peak. If \(\mathrm{E}_1,\mathrm{E}_2\) are peaks, then \(\mathrm{E}_1\cup \mathrm{E}_2\) is a peak. (Let \(f_{1},f_{2}\in \mathrm{A}\) with \(f_{i} = 1\) on \(\mathrm{E}_i\), \(|f_i| < 1\) on \(\mathrm{X - E}_i\); let \(g_{i} = \frac{1}{4} (1 - f_{i})^{1 / 3}\in \mathrm{A}\); then \(1 - g_{1}g_{2} = 1\) on \(\mathrm{E}_1\cup \mathrm{E}_2\), \(|1 - g_1g_2| < 1\) on \(\mathrm{X - (E_1\cup E_2)}\).)

\begin{enumerate}
\def\labelenumi{\alph{enumi})}
\setcounter{enumi}{4}
\item
  Let E be an intersection of peaks. Let I be the set of \(f \in A\) vanishing on E. Then \(f \mapsto f|E\) defines an isometry from A/I onto the algebra A\textbar E of restrictions to E of functions in A.
\item
  Let E be a peak. Let \(E' \subset E\) be a peak relative to A\textbar E. Then \(E'\) is a peak relative to A.
\item
  Every antisymmetric subset of X is an intersection of peaks. (Use d) and f).)
\item
  For every antisymmetric subset K of X, the algebra A\textbar K is a closed subalgebra of \(\mathcal{C}(\mathrm{K})\). (Use \(e\) and \(g\).)
\item
  Let R be the set of real parts of elements of A. Suppose that \(X \in \mathfrak{R}\) and that R is stable under multiplication. Then X consists of a single point. (Let \(x_{0} \in X\). For \(u \in R\), there exists a unique \(f \in A\) such that \(u = \mathcal{R}(f)\) and \(f(x_{0}) \in \mathbf{R}\); set \(\mathrm{N}(u) = \|f\|\). Then R is a real Banach space for N; the closed graph theorem proves that multiplication in R is separately continuous, hence continuous. Deduce from this a Banach algebra structure on \(S = R + iR\). Let \(p \in R\) with \(p(x) > 0\) for all \(x \in X\). Show that, for every character \(\chi\) of S, one has \(\mathcal{R}(\chi(p)) > 0\), and deduce that \(\log(p) \in R\). Assuming \(X \neq \{x_{0}\}\), one can construct an \(f \in A\) such that \(0 < \mathcal{R}(f) \leqslant 1\) on X, \(f(x_{0}) \in \mathbf{R}\), and such that \(\|f\|\) is arbitrarily large. According to what precedes, there exists \(V \in A\) with \(|V|^{2} = \mathcal{R}(f)\). We have \(\|V\| \leqslant 1\), hence \(\mathrm{N}(\mathcal{R}(V)) \leqslant 2\), but
\end{enumerate}

\[
\mathrm{N} (\mathcal {R} (\mathrm{V}) ^ {2}) \geqslant \frac {1}{2} \left(\| \mathrm{V} ^ {2} + f \| - | \mathcal {I} (\mathrm{V}) (x _ {0}) ^ {2} |\right)
\]

is arbitrarily large.)

\begin{enumerate}
\def\labelenumi{\alph{enumi})}
\setcounter{enumi}{9}
\tightlist
\item
  If R is stable under multiplication, then A = C(X). (Use a), c), h), i).) Consequently, if A \(\neq\) C(X), there exists \(u \in R\) such that \(u^{2} \notin R\).
\end{enumerate}

\begin{enumerate}
\def\labelenumi{\arabic{enumi})}
\setcounter{enumi}{10}
\tightlist
\item
  Let X be a metric space, whose distance is denoted by d. A function \(f: X \to C\) is said to be Lipschitz if there exists a real number \(c \geqslant 0\) such that, for all \(x\) and \(y\) in X, one has
\end{enumerate}

\[
| f (x) - f (y) | \leqslant c d (x, y)
\]

(cf.~FVR, IV, p.~7, def. 1). We denote by \(\mathrm{Lip}_b(\mathrm{X})\) the set of bounded Lipschitz functions on X.

\begin{enumerate}
\def\labelenumi{\alph{enumi})}
\item
  The set \(\mathrm{Lip}_b(\mathrm{X})\) is a subalgebra of \(\mathcal{C}(\mathrm{X})\).
\item
  Equipped with the norm
\end{enumerate}

\[
\| f\| = \sup_{x\in \mathrm{X}}|f(x)| + \sup_{\substack{x,y\in \mathrm{X}\\ x\neq y}}\frac{|f(x) - f(y)|}{d(x,y)},
\]

and with the involution \(f \mapsto \overline{f}\), the algebra \(\mathrm{Lip}_b(\mathrm{X})\) is a commutative involutive Banach algebra. It is not a stellar algebra in general.

\begin{enumerate}
\def\labelenumi{\alph{enumi})}
\setcounter{enumi}{2}
\item
  For any commutative Banach algebra A, we denote by \(\mathsf{X}_{m}(\mathsf{A})\) the metric space \(\mathsf{X}(\mathsf{A})\) endowed with the distance induced by the norm of the dual \(A'\) of A. The metric space \(\mathsf{X}_{m}(\mathsf{A})\) is complete and of diameter \(\leqslant 2\) .
\item
  Let A be a unital commutative Banach algebra. The Gelfand transform maps A into \(\operatorname{Lip}_{b}(\mathsf{X}_{m}(\mathsf{A}))\). If A is semisimple, then it induces a homeomorphism from A onto its image.
\end{enumerate}

In the sequel, we assume that X is a compact metric space. We denote by \(A = \operatorname{Lip}_{b}(X)\).

\begin{enumerate}
\def\labelenumi{\alph{enumi})}
\setcounter{enumi}{4}
\item
  The map \(\varphi\) which sends \(x \in X\) to the character \(f \mapsto f(x)\) is a homeomorphism from \(X\) onto \(X(A)\).
\item
  Let \(d_{1}\) be the distance on X induced by the homeomorphism \(\varphi\). Then \(d_{1}\) is equivalent to d.
\end{enumerate}

